\begin{abstract}
\addchaptertocentry{\abstractname} % Agregar resumen al índice
Se desarrolló una metodología para estimar la Potencia Límite de Canal de la Central Nuclear Atucha II, basada en redes neuronales, con tiempos de cómputo órdenes de magnitud inferiores a los códigos termohidráulicos empleados en Nucleoeléctrica Argentina S.A. Esto permite disminuir considerablemente los tiempos y costos computacionales de los cálculos termohidráulicos para las estrategias de recambio de combustibles. 

Adicionalmente, se propone una metodología novedosa para estimar el Flujo Crítico de Calor, que incorpora redes neuronales informadas por física y redes neuronales de multifidelidad. El enfoque integra datos provenientes de ensayos experimentales con geometrías simples y con la geometría del \textit{bundle} de la Central Nuclear Atucha II, y ensayos numéricos. 

Para el desarrollo del trabajo se aplicaron conocimientos en redes neuronales, redes neuronales informadas por física, redes neuronales de multifidelidad, y su entrenamiento y optimización.


% Dentro del diseño termohidráulico de reactores nucleares, la Potencia Límite de Canal establece el límite al cual un canal refrigerante puede operar de manera segura. Esto lo convierte en un parámetro fundamental para el cálculo de estrategias de recambio de combustibles, y una alta eficiencia computacional y precisión en su estimación son cruciales para garantizar la operación segura de la central.

% Este trabajo evalúa dos posibles mejoras en ambas categorías. En primer lugar, se presenta una metodología para la estimación de la Potencia Límite de Canal con redes neuronales, entrenadas con datos de simulaciones de código de subcanales, que permiten la evaluación de estados del reactor en instantes. En segundo lugar, se propone un método novedoso para estimar el Flujo Crítico de Calor, un parámetro fundamental para la estimación de la Potencia Límite de Canal, a partir de datos de variada fidelidad: datos experimentales de ensayos con la geometría específica del combustible de la Central Nuclear Atucha II, simulaciones de códigos de subcanales con la misma geometría y las metodologías vigentes de estimación de Flujo Crítico de Calor, y datos experimentales de geometrías sencillas. Esta metodología se combina con Redes Neuronales Informadas por Física, para asegurar predicciones físicamente válidas, lo cual es fundamental en la industria nuclear.
\end{abstract}